\subsection{3회차: 이득 항등식과 자유 작용의 한계 (증명 노트)}
{\footnotesize\noindent 원본: \path{rounds/R3_G_klartag/notes.md}\par}
\VerbatimInput[breaklines=true,breakanywhere=true,fontsize=\scriptsize,formatcom=\notesfont]{notes_tmp/R3_G_klartag_notes.md}

\subsection{3회차: 비자유 작용, 텐서 안정성, no-go (증명 노트)}
{\footnotesize\noindent 원본: \path{rounds/R3_non_free/notes.md}\par}
\VerbatimInput[breaklines=true,breakanywhere=true,fontsize=\scriptsize,formatcom=\notesfont]{notes_tmp/R3_non_free_notes.md}

\subsection{3회차: 과정의 한계, Lemma N (증명 노트)}
{\footnotesize\noindent 원본: \path{rounds/R3_process/notes.md}\par}
\VerbatimInput[breaklines=true,breakanywhere=true,fontsize=\scriptsize,formatcom=\notesfont]{notes_tmp/R3_process_notes.md}

\subsection{4회차: 정확한 $\chi^2$ 법칙 등 (심사자 노트)}
{\footnotesize\noindent 원본: \path{rounds/R4_perceptron_theory/referee/notes.md}\par}
\VerbatimInput[breaklines=true,breakanywhere=true,fontsize=\scriptsize,formatcom=\notesfont]{notes_tmp/R4_perceptron_theory_referee_notes.md}

\subsection{5회차: Leech 격자의 계산 (노트)}
{\footnotesize\noindent 원본: \path{rounds/R5_leech13/notes.md}\par}
\VerbatimInput[breaklines=true,breakanywhere=true,fontsize=\scriptsize,formatcom=\notesfont]{notes_tmp/R5_leech13_notes.md}

\subsection{5회차: 일반 $p$ 의 보조정리 (증명 노트)}
{\footnotesize\noindent 원본: \path{rounds/R5_sl2p_theory/notes.md}\par}
\VerbatimInput[breaklines=true,breakanywhere=true,fontsize=\scriptsize,formatcom=\notesfont]{notes_tmp/R5_sl2p_theory_notes.md}

\subsection{7회차: 원분 접착 (노트)}
{\footnotesize\noindent 원본: \path{rounds/R7_gluing/notes.md}\par}
\VerbatimInput[breaklines=true,breakanywhere=true,fontsize=\scriptsize,formatcom=\notesfont]{notes_tmp/R7_gluing_notes.md}

\subsection{10회차: Weil 격자 --- 빌더 A 노트}
{\footnotesize\noindent 원본: \path{rounds/R10_weil_lattices/A/notes.md}\par}
\VerbatimInput[breaklines=true,breakanywhere=true,fontsize=\scriptsize,formatcom=\notesfont]{notes_tmp/R10_weil_lattices_A_notes.md}

\subsection{10회차: Weil 격자 --- 빌더 B 노트}
{\footnotesize\noindent 원본: \path{rounds/R10_weil_lattices/B/notes.md}\par}
\VerbatimInput[breaklines=true,breakanywhere=true,fontsize=\scriptsize,formatcom=\notesfont]{notes_tmp/R10_weil_lattices_B_notes.md}

\subsection{11회차: Mordell--Weil 재정식화의 증명 (T1 노트)}
{\footnotesize\noindent 원본: \path{rounds/R11_mordell_weil/T1/notes.md}\par}
\VerbatimInput[breaklines=true,breakanywhere=true,fontsize=\scriptsize,formatcom=\notesfont]{notes_tmp/R11_mordell_weil_T1_notes.md}

\subsection{11회차: 조건부 주장에 대한 반박 시도 (T2 노트)}
{\footnotesize\noindent 원본: \path{rounds/R11_mordell_weil/T2/notes.md}\par}
\VerbatimInput[breaklines=true,breakanywhere=true,fontsize=\scriptsize,formatcom=\notesfont]{notes_tmp/R11_mordell_weil_T2_notes.md}

\subsection{12회차: $p=61$ 의 체계적 시험 (노트)}
{\footnotesize\noindent 원본: \path{rounds/R12_p61/notes.md}\par}
\VerbatimInput[breaklines=true,breakanywhere=true,fontsize=\scriptsize,formatcom=\notesfont]{notes_tmp/R12_p61_notes.md}

\subsection{12회차: 독립 검증 노트}
{\footnotesize\noindent 원본: \path{rounds/R12_p61/verify/notes.md}\par}
\VerbatimInput[breaklines=true,breakanywhere=true,fontsize=\scriptsize,formatcom=\notesfont]{notes_tmp/R12_p61_verify_notes.md}

\subsection{14회차: 문헌 점검 보고}
{\footnotesize\noindent 원본: \path{rounds/R14_leech_tensor/novelty_check.md}\par}
\VerbatimInput[breaklines=true,breakanywhere=true,fontsize=\scriptsize,formatcom=\notesfont]{notes_tmp/R14_leech_tensor_novelty_check.md}

\subsection{14회차: 심사 보고 A}
{\footnotesize\noindent 원본: \path{rounds/R14_leech_tensor/referee_A.md}\par}
\VerbatimInput[breaklines=true,breakanywhere=true,fontsize=\scriptsize,formatcom=\notesfont]{notes_tmp/R14_leech_tensor_referee_A.md}

\subsection{14회차: 심사 보고 B}
{\footnotesize\noindent 원본: \path{rounds/R14_leech_tensor/referee_B.md}\par}
\VerbatimInput[breaklines=true,breakanywhere=true,fontsize=\scriptsize,formatcom=\notesfont]{notes_tmp/R14_leech_tensor_referee_B.md}

\subsection{15회차: 심사 보고 A}
{\footnotesize\noindent 원본: \path{rounds/R15_rogers/referee.md}\par}
\VerbatimInput[breaklines=true,breakanywhere=true,fontsize=\scriptsize,formatcom=\notesfont]{notes_tmp/R15_rogers_referee.md}

\subsection{15회차: 심사 보고 B}
{\footnotesize\noindent 원본: \path{rounds/R15_rogers/referee_B.md}\par}
\VerbatimInput[breaklines=true,breakanywhere=true,fontsize=\scriptsize,formatcom=\notesfont]{notes_tmp/R15_rogers_referee_B.md}

\subsection{16회차: 몫곡선 사전의 심사 보고}
{\footnotesize\noindent 원본: \path{rounds/R16_quotient_curve/referee.md}\par}
\VerbatimInput[breaklines=true,breakanywhere=true,fontsize=\scriptsize,formatcom=\notesfont]{notes_tmp/R16_quotient_curve_referee.md}

\subsection{17회차: Gargava 정리 16 의 점검 노트}
{\footnotesize\noindent 원본: \path{rounds/R17_gargava_check/notes.md}\par}
\VerbatimInput[breaklines=true,breakanywhere=true,fontsize=\scriptsize,formatcom=\notesfont]{notes_tmp/R17_gargava_check_notes.md}

\subsection{17회차: 심사 보고}
{\footnotesize\noindent 원본: \path{rounds/R17_gargava_check/referee.md}\par}
\VerbatimInput[breaklines=true,breakanywhere=true,fontsize=\scriptsize,formatcom=\notesfont]{notes_tmp/R17_gargava_check_referee.md}

\subsection{18회차: 14·15회차 결합 점검 보고 (통합 진술 포함)}
{\footnotesize\noindent 원본: \path{rounds/R18_integration/report.md}\par}
\VerbatimInput[breaklines=true,breakanywhere=true,fontsize=\scriptsize,formatcom=\notesfont]{notes_tmp/R18_integration_report.md}

